\documentclass[11pt,a4paper]{article}
\usepackage[utf8]{inputenc}
\usepackage[T1]{fontenc}
\usepackage{amsmath,amssymb,amsthm}
\usepackage{geometry}
\usepackage{hyperref}
\usepackage{fancyhdr}
\usepackage{titlesec}
\usepackage{enumitem}
\usepackage{array}
\usepackage{longtable}
\usepackage{booktabs}
\usepackage{mathrsfs}
\usepackage{bbm}

\geometry{margin=1in}

\pagestyle{fancy}
\fancyhf{}
\fancyhead[C]{\small Arthur Cayley --- Collected Mathematical Papers, Vol.~XI}
\fancyfoot[C]{\thepage}
\renewcommand{\headrulewidth}{0.4pt}

\theoremstyle{definition}
\newtheorem*{remark}{Remark}

\title{\textbf{The Collected Mathematical Papers of Arthur Cayley}\\[6pt]
\large Volume XI, Pages 351--400\\[4pt]
\normalsize Papers 766--774}
\author{}
\date{}

\begin{document}

\maketitle
\thispagestyle{empty}

\tableofcontents
\newpage

%===========================================================================
\section[On the Geodesic Curvature of a Curve on a Surface]{766. On the Geodesic Curvature of a Curve on a Surface.\footnote{From the \textit{Proceedings of the London Mathematical Society}, vol.~xii. (1881), pp.~129--132. Read March 10, 1881.}}
%===========================================================================

\noindent\textsc{[Continued from previous pages.]}

\subsection*{The Special Curves, $p = \text{constant}$ and $q = \text{constant}$}

Consider the curve $p = \text{const}$. For this curve $p' = 0$, $p'' = 0$; therefore also $Gq'^2 = 1$, and, if $R$ be the radius of geodesic curvature, then
\[
\frac{1}{R} = \frac{\sqrt{G_1}}{G}.
\]
Similarly for the curve $q = \text{const}$. Here $q' = 0$, $q'' = 0$; therefore $Ep'^2 = 1$, and, if $S$ be the radius of geodesic curvature, then
\[
\frac{1}{S} = \frac{\sqrt{E_2}}{E}.
\]
These values of $R$ and $S$ are interesting for their own sakes, and they will be introduced into the expression for the radius of geodesic curvature $\rho$ of the general curve.

\subsection*{Transformed Formula for the Radius of Geodesic Curvature}

From the values of $\dfrac{1}{R}$, $\dfrac{1}{S}$ we have
\[
\frac{1}{\rho} = \sqrt{\frac{E}{G}} \frac{p'}{q'} + \sqrt{\frac{G}{E}} \frac{q'}{p'} + \sqrt{EG} \left\{ p'q'' - p''q' + \frac{1}{2}(Lp' + Mq') \right\},
\]
where the term in $\{\ \}$ is
\[
= 2\mathfrak{A}p' - \tfrac{1}{2}E_2 p'^3 + \mathfrak{B}p''q' + (\mathfrak{C}p'q'' + \mathfrak{D}q'^2) - \tfrac{1}{2}G_1 q'^3.
\]
The terms in $\mathfrak{A}$ are
\[
= -\tfrac{1}{2}p'(1 - Ep'^2) = -\tfrac{1}{2}p'(2Fp'q' + Gq'^2),
\]
and those in $\mathfrak{D}$ are
\[
= -\tfrac{1}{2}q'(1 - Gq'^2) = -\tfrac{1}{2}q'(Ep'^2 + 2Fp'q').
\]
Hence the whole expression contains the factor $p'q'$, and is, in fact,
\[
= p'q'\{Lp' + Mq'\};
\]
and the formula thus is
\[
\frac{1}{\rho} = \frac{\sin^2\phi}{R} + \frac{\sin^2\psi}{S} + \frac{\sin\omega \cdot p'q'}{\sqrt{EG - F^2}}\{Lp' + Mq'\},
\]
where $L$ and $M$ are certain expressions involving the coefficients $E$, $F$, $G$ and their derivatives.

Taking $\phi$, $\psi$ to be the inclination of the curve to the curves $q = \text{const.}$, $p = \text{const.}$, respectively, and $\omega\,(= \phi + \psi)$ the inclination of these two curves to each other, then
\[
\cos\phi = \frac{Fp' + Gq'}{\sqrt{G}},\quad \sin\phi = \sqrt{E}\,p'\sqrt{EG - F^2},\quad \text{etc.}
\]
hence $\sin\phi = p'\sqrt{E}\sin\omega$, $\sin\psi = q'\sqrt{G}\sin\omega$, and the formula may also be written
\[
\frac{1}{\rho} = \frac{\sin^2\phi}{R \sin^2\omega} + \frac{\sin^2\psi}{S \sin^2\omega} + \frac{\sin\phi \sin\psi}{\sqrt{EG - F^2}}\{Lp' + Mq'\}.
\]

\subsection*{The Orthotomic Case $F = 0$, or $ds^2 = Edp^2 + Gdq^2$}

The formula becomes in this case much more simple. We have
\[
1 = Ep'^2 + Gq'^2,\quad \sqrt{EG - F^2} = \sqrt{EG},\quad \omega = 90^\circ,\quad \sin\phi = \cos\psi;
\]
and the term $Lp' + Mq'$ becomes $= EG_1 - E_1G$, if, as before, $E_1$, $G_1$ denote the complete differential coefficients $E_1p' + E_2q'$ and $G_1p' + G_2q'$. The formula then is
\[
\frac{1}{\rho} - \frac{\cos^2\phi}{R} - \frac{\sin^2\phi}{S} = \sqrt{EG}(p'q'' - p''q') + \tfrac{1}{2}p'q'\sqrt{EG}(EG_1 - E_1G),
\]
where the values $\dfrac{1}{R}$ and $\dfrac{1}{S}$ are now $= \dfrac{\sqrt{G_1}}{G}$ and $\dfrac{\sqrt{E_2}}{E}$ respectively. But we have moreover $\phi = \tan^{-1}\dfrac{\sqrt{G}\,q'}{\sqrt{E}\,p'}$, and thence
\[
\phi' = \sqrt{EG}(p'q'' - p''q') - \tfrac{1}{2}\phi'\sqrt{EG}(EG_1 - E_1G);
\]
or the formula finally is
\[
\frac{1}{\rho} - \frac{\cos^2\phi}{R} - \frac{\sin^2\phi}{S} = \phi',
\]
which is Liouville's formula referred to at the beginning of the present paper. It will be recollected that $\phi'$ is the differential coefficient with respect to the arc $s$ of the curve.

\subsection*{Addition}

Since the foregoing paper was written, I have succeeded in obtaining a like interpretation of the term
\[
\sqrt{EG - F^2}(p'q'' - p''q') + \tfrac{1}{2}p'q'(Lp' + Mq'),
\]
which belongs to the general case. I find that these terms are, in fact, $= -\phi' + \omega_1 p'$; or, what is the same thing (since $\omega = \psi + \phi$ and therefore $\omega_1 p' + \omega_2 q' = \psi' + \phi'$), are $= \psi' - \omega_1$. It will be recollected that $\psi$ is the inclination of the curve to the curve $q = c$, which passes through a given point of the curve, $\psi'$ is the variation of $\psi$ corresponding to the passage to the consecutive point of the curve, viz., $\psi + \psi'\,ds$ is the inclination at this consecutive point to the curve $q = c + dc$, which passes through the consecutive point; $\omega$ is the inclination to each other of the curves $p = b$, $q = c$, which pass through the given point of the curve, $\omega_1$ the variation corresponding to the passage along the curve $q = c$, viz., $\omega + \omega_1\,ds$ is the inclination to each other of the curves $p = b + db$, $q = c$; and the like as regards $\phi$ and $\omega_2$.

For the demonstration, we have, as above,
\[
\cos\psi = \frac{\sqrt{G}\,q'}{\sqrt{EG - F^2}},\quad \sin\psi = \frac{\sqrt{G}\,p'}{\sqrt{EG - F^2}},
\]
where $\sqrt{EG - F^2} = \sqrt{EG - F^2}$; and moreover $Ep'^2 + 2Fp'q' + Gq'^2 = 1$. In virtue of this last equation,
\[
\sqrt{EG - F^2}\,\psi' + (Fp' + Gq')^2 = G;
\]
and we have
\[
\psi' = -\sqrt{EG - F^2}(p'q'' - p''q') + \frac{D}{(EG - F^2)^{3/2}},
\]
where
\[
D = (Fp' + Gq')p'\dot{V} - Vp'(Fp' + Gq')^\cdot;
\]
or, since $V^2 = EG - F^2$, and thence $2V\dot{V} = \dot{E}G + E\dot{G} - 2F\dot{F}$, we have
\[
D = \tfrac{1}{2}\dot{V}\{(Fp' + Gq')(\dot{E}G + E\dot{G} - 2F\dot{F}) - 2(EG - F^2)(Fp' + Gq')^\cdot\}.
\]
Substituting herein for $\dot{E}$, $\dot{F}$, $\dot{G}$ their values $E_1p' + E_2q'$, $F_1p' + F_2q'$, $G_1p' + G_2q'$, the term in $\{\ \}$ becomes
\[
= Ip'^2 + Jp'q' + Kq'^2,
\]
where
\[
\begin{aligned}
I &= FGE_1 - 2EGF_1 + EFG_1,\\
J &= G'E_1 - 2FGF_1 + (-EG + 2F^2)G_1 + FGE_2 - 2EGF_2 + EFG_2,\\
K &= G'E_2 - 2FGF_2 + (-EG + 2F^2)G_2.
\end{aligned}
\]
But from the equation $\omega = \tan^{-1}\dfrac{\sqrt{EG - F^2}}{F}$, differentiating in regard to $p$, we obtain
\[
\omega_1 = \frac{FGE_1 - 2EGF_1 + EFG_1}{(EG - F^2)^{3/2}} = \frac{I}{(EG - F^2)^{3/2}};
\]
or, for $p$ writing $p + p'(\dot{E}p'^2 + 2\dot{F}p'q' + \dot{G}q'^2)\sqrt{EG - F^2} = \dot{E}p' + \dotsb$, we have
\[
\psi' - \omega_1 p' = -\sqrt{EG - F^2}(p'q'' - p''q') + (Ip'^2 + Jp'q' + Kq'^2).
\]
The terms in $p''$ destroy each other, and the form thus is
\[
\psi' - \omega_1 p' = -\sqrt{EG - F^2}(p'q'' - p''q') - \tfrac{1}{2}p'q'(Lp' + Mq'),
\]
where
\[
L = \frac{K}{p'^2} - \frac{I}{q'^2},\quad M = \frac{K}{p'q'} - \frac{J}{p'^2},
\]
and, upon substituting herein for $I$, $J$, $K$ their values, we find
\[
\begin{aligned}
L &= -GE_1 + EG_1 + \frac{2F(FG_1 - GF_1) - F(EG_1 - GE_1) + 2E(FF_1 - EF_1)}{EG - F^2},\\
M &= -GE_2 + EG_2 + \frac{2F(FG_2 - GF_2) - F(EG_2 - GE_2) + 2E(FF_2 - EF_2)}{EG - F^2};
\end{aligned}
\]
viz., these are the values denoted above by the same letters $L$, $M$. The final result thus is
\[
\frac{1}{\rho} - \frac{\sin^2\phi}{R} - \frac{\sin^2\psi}{S} = \psi' - \omega_1 p' = \phi' - \omega_2 q',
\]
where the meanings of the symbols have been already explained. A formula substantially equivalent to this, but in a different (and scarcely properly explained) notation, is given, Aoust, ``Th\'eorie des coordonn\'ees curvilignes quelconques,'' \textit{Annali di Matem.}, t.~ii. (1868), pp.~39--64; and I was, in fact, led thereby to the foregoing further investigation.

As to the definition of the radius of geodesic curvature, I remark that, for a curve on a given surface, if $PQ$ be an infinitesimal arc of the curve, then if from $Q$ we let fall the perpendicular $QM$ on the tangent plane at $P$ (the point $M$ being thus a point on the tangent $PT$ of the curve), and if from $M$, in the tangent plane and at right angles to the tangent, we draw $MN$ to meet the osculating plane of the curve in $N$, then $MN$ is in fact equal to the infinitesimal arc $QQ'$ mentioned near the beginning of the present paper, and the radius of geodesic curvature $\rho$ is thus a length such that $2\rho \cdot MN = PQ^2$.

\newpage
%===========================================================================
\section[On the Gaussian Theory of Surfaces]{767. On the Gaussian Theory of Surfaces.\footnote{From the \textit{Proceedings of the London Mathematical Society}, vol.~xii. (1881), pp.~187--192. Read June 9, 1881.}}
%===========================================================================

In the Memoir, Bour, ``Th\'eorie de la d\'eformation des surfaces'' (\textit{Jour.~de l'\'Ec.~Polyt.}, Cah.~39 (1862), pp.~1--148), the author, working with the form $ds^2 = dv^2 + g^2 du^2$ as a special case of Gauss's formula $ds^2 = Edp^2 + 2Fdpdq + Gdq^2$, obtains (p.~29) the following equations which he calls fundamental:
\[
\begin{aligned}
\frac{dH}{dv} - \frac{dT}{du} &= TV \frac{dg}{g\,dv},\\
\frac{dT}{dv} - \frac{dH_1}{du} &= -\frac{TV}{g} \frac{dg}{dv},
\end{aligned}
\tag{IV.}
\]
where $T$, $H$ is written to denote $\dfrac{d\cdot}{du}$, $\dfrac{d\cdot}{dv}$, and where (see p.~26)
\begin{itemize}[noitemsep]
\item $H$ is the curvature of the normal section containing the tangent to the curve $v = \text{constant}$,
\item $H_1$ is the curvature of the normal section at right angles to the preceding, containing the tangent to the (geodesic) curve $u = \text{constant}$,
\item $T$ is the torsion of the same geodesic curve;
\end{itemize}
or, what is the same thing (see p.~25), the quadric equation for the determination of the principal radii of curvature at the point of the surface is
\[
\left(\frac{1}{\rho} - H\right)\left(\frac{1}{\rho} - H_1\right) - T^2 = 0.
\]

Writing for greater convenience $K$ in place of the suffixed letter $H_1$, also $V$ instead of $g$, so that the differential formula is $ds^2 = dv^2 + V^2 du^2$, the equations become
\[
\begin{aligned}
\frac{dH}{dv} - \frac{dT}{du} &= TV \frac{V_1}{V},\\
\frac{dT}{dv} - \frac{dK}{du} &= -\frac{TV}{V} V_1;
\end{aligned}
\]
or, if we use the suffix $1$ to denote differentiation in regard to $v$, and the suffix $2$ to denote differentiation in regard to $u$, then the equations are
\[
\begin{aligned}
H_1 - T_2 &= \frac{TV V_1}{V},\\
T_1 - K_2 &= -\frac{TV V_1}{V};
\end{aligned}
\]
or, what is the same thing,
\[
\begin{aligned}
T_1 + H_1 V + (H - K)V_1 &= 0,\\
T_1 V + 2TV_1 + K_2 &= 0.
\end{aligned}
\]

I wish to show how these formulae connect themselves with formulae belonging to the general form $ds^2 = Edp^2 + 2Fdpdq + Gdq^2$. These involve not only Gauss's coefficients $E$, $F$, $G$, but also the coefficients $E'$, $F'$, $G'$ belonging to the inflexional tangents; and, for convenience, I quote the system of definitions, Salmon's \textit{Geometry of Three Dimensions}, 3rd ed., 1874, p.~251, viz.
\[
\begin{aligned}
dx,\; dy,\; dz &= a\,dp + a'\,dq,\; b\,dp + b'\,dq,\; c\,dp + c'\,dq;\\
d^2x &= \alpha\,dp^2 + 2\alpha'\,dp\,dq + \alpha''\,dq^2,\\
d^2y &= \beta\,dp^2 + 2\beta'\,dp\,dq + \beta''\,dq^2,\\
d^2z &= \gamma\,dp^2 + 2\gamma'\,dp\,dq + \gamma''\,dq^2;\\
A,\; B,\; C &= bc' - b'c,\; ca' - c'a,\; ab' - a'b;\quad V^2 = EG - F^2;\\
E' &= A\alpha + B\beta + C\gamma,\quad F' = A\alpha' + B\beta' + C\gamma',\quad G' = A\alpha'' + B\beta'' + C\gamma'',
\end{aligned}
\]
so that $E'$, $F'$, $G'$ are, in fact, the determinants
\[
\begin{vmatrix}
a & b & c \\
a' & b' & c' \\
\alpha & \beta & \gamma
\end{vmatrix},\quad
\begin{vmatrix}
a & b & c \\
a' & b' & c' \\
\alpha' & \beta' & \gamma'
\end{vmatrix},\quad
\begin{vmatrix}
a & b & c \\
a' & b' & c' \\
\alpha'' & \beta'' & \gamma''
\end{vmatrix}.
\]
The equation for the determination of the principal radii of curvature is
\[
(E'\rho - EV)(G'\rho - GV) - (F'\rho - FV)^2 = 0.
\]

which, in the particular case $F = 0$ (and therefore $V^2 = EG$), becomes
\[
(E'\rho - EV)(G'\rho - GV) - F'^2\rho^2 = 0,
\]
or, as this may be written,
\[
\left(\frac{1}{\rho} - \frac{E'}{EV}\right)\left(\frac{1}{\rho} - \frac{G'}{GV}\right) - \frac{F'^2}{EGV^2} = 0,
\]
an equation which corresponds with Bour's form
\[
\left(\frac{1}{\rho} - H\right)\left(\frac{1}{\rho} - K\right) - T^2 = 0,
\]
and becomes identical with it, if
\[
E' = EVK,\quad G' = GVH,\quad F' = -V^2T.
\]
But, making $p$, $q$ correspond to Bour's variables, $p$ to $v$, and $q$ to $u$, it is necessary to show that the foregoing values (and not the interchanged values $E' = GVH$, $G' = EVK$) are the correct ones. We have, Salmon, p.~254,
\[
\begin{vmatrix}
dq, & pE' - VE, & pF' - VF \\
-dp, & pF' - VF, & pG' - VG
\end{vmatrix}
= 0;
\]
or, putting herein $F = 0$, the equations may be written
\[
\frac{dq}{-dp} = \frac{E'\left(\dfrac{1}{\rho} - \dfrac{VE}{E'}\right)}{F'\left(\dfrac{1}{\rho} - \dfrac{VF}{F'}\right)} = \frac{F'\left(\dfrac{1}{\rho} - \dfrac{VF}{F'}\right)}{G'\left(\dfrac{1}{\rho} - \dfrac{VG}{G'}\right)};
\]
or, we see that to $dq = 0$ corresponds the value $\dfrac{1}{\rho} = \dfrac{VE}{E'}$, and to $dp = 0$ the value $\dfrac{1}{\rho} = \dfrac{VG}{G'}$. Hence the former of these values of $\dfrac{1}{\rho}$ corresponds to Bour's $du = 0$, that is, to his $\dfrac{1}{\rho} = K$; and the latter to Bour's $dv = 0$, that is, to his $\dfrac{1}{\rho} = H$; or the values are, as stated,
\[
E' = EVK,\quad G' = GVH.
\]
The formula $ds^2 = Edp^2 + 2Fdpdq + Gdq^2$ agrees with Bour's $ds^2 = dv^2 + g^2 du^2$, if $p = u$, $q = v$, $E = 1$, $F = 0$, $G = g^2$. With these values, $V^2 = EG - F^2 = g^2$, or say $g = V$, and Bour's equation is, as it was before written, $ds^2 = dv^2 + V^2 du^2$. And we have to find the three equations which, putting therein $p = u$, $q = v$, $E = 1$, $F = 0$, $G = V^2$, $E' = VK$, $F' = -VT$, $G' = VH$, reduce themselves to Bour's equations.

The first of these is nothing else than the equation for the measure of curvature, viz. Salmon, p.~262 (but, using the suffix $1$ for $v$ and $2$ for $u$),
\[
4V^2\{HK - T^2\} = (2VV_{12})^2 - 2V_2(2VV_{11} + 2V_1^2) + 2V_1(2VV_{22} + 2V_2^2).
\]
In fact, writing herein $E = 1$, $F = 0$, and therefore the differential coefficients of $E$ and $F$ each $= 0$, the equation becomes
\[
4V^4\{HK - T^2\} = (2VV_{12})^2 - 2V_2(2VV_{11} + 2V_1^2) + 2V_1(2VV_{22} + 2V_2^2),
\]
which is
\[
4V^4\{HK - T^2\} = (2VV_{12})^2 - 2V_2(2VV_{11} + 2V_1^2) + 2V_1(2VV_{22} + 2V_2^2) = -4V^4V_{22};
\]
or finally it is
\[
V_{22} = V(T^2 - HK).
\]

The other two of Bour's equations are derived from equations which give respectively the values of $E_2' - F_1'$ and $F_2' - G_1'$; viz. starting from the equations
\[
\begin{aligned}
E' &= A\alpha + B\beta + C\gamma,\\
F' &= A\alpha' + B\beta' + C\gamma',\\
G' &= A\alpha'' + B\beta'' + C\gamma'',
\end{aligned}
\]
we see at once that $E_2'$ and $F_1'$ contain, $E_2'$ the terms $A\alpha_2 + B\beta_2 + C\gamma_2$, and $F_1'$ the terms $A\alpha_1' + B\beta_1' + C\gamma_1'$, which are equal to each other ($\alpha_2 = \alpha_1'$ since $\alpha$ and $\alpha'$ are the differential coefficients $x_1$, $x_2$ of $x$, and so $\beta_2 = \beta_1'$ and $\gamma_2 = \gamma_1'$). Hence
\[
E_2' - F_1' = A_2\alpha + B_2\beta + C_2\gamma - A_1\alpha' - B_1\beta' - C_1\gamma',
\]
and similarly
\[
F_2' - G_1' = A_2\alpha' + B_2\beta' + C_2\gamma' - A_1\alpha'' - B_1\beta'' - C_1\gamma''.
\]

Here, from the values of $A$, $B$, $C$, we have
\[
\begin{aligned}
A &= bc' - cb'; & A_1 &= \beta c' - \gamma b' + b\gamma' - c\beta'; & A_2 &= \beta'c' - \gamma'b' + b\gamma'' - c\beta'';\\
B &= ca' - ac'; & B_1 &= \gamma a' - \alpha c' + c\alpha' - a\gamma'; & B_2 &= \gamma'a' - \alpha'a' + c\alpha'' - a\gamma'';\\
C &= ab' - ba'; & C_1 &= \alpha b' - \beta a' + a\beta' - b\alpha'; & C_2 &= \alpha'b' - \beta'a' + a\beta'' - b\alpha'';
\end{aligned}
\]
and, substituting, we find
\[
\begin{aligned}
E_2' - F_1' &= 2\alpha'\alpha\alpha' + \alpha\alpha''\alpha - 2\alpha\alpha'\alpha'' - \alpha'\alpha''\alpha,\\
F_2' - G_1' &= \dotsb
\end{aligned}
\]
if, for shortness, $\alpha'\alpha\alpha'$ denotes the determinant
\[
\begin{vmatrix}
\alpha' & \alpha & \alpha \\
\beta' & \beta & \beta \\
\gamma' & \gamma & \gamma'
\end{vmatrix}
\]
and so for the other like symbols. Observe that, with
\[
\alpha,\; \alpha',\; \alpha,\; \alpha',\; \alpha''\\
b,\; b',\; \beta,\; \beta',\; \beta''\\
c,\; c',\; \gamma,\; \gamma',\; \gamma''
\]
we have in all 10 determinants, viz. these are $\alpha\alpha'\alpha = E'$; $\alpha\alpha'\alpha' = F'$; $\alpha\alpha'\alpha'' = G'$; $\alpha\alpha'\alpha''$; and the six determinants $\alpha\alpha\alpha'$, $\alpha\alpha'\alpha''$, $\alpha\alpha''\alpha$; $\alpha'\alpha\alpha$, $\alpha'\alpha'\alpha''$, $\alpha'\alpha''\alpha$.

The foregoing expressions of $E_2' - F_1'$ and $F_2' - G_1'$ respectively, substituting therein for the determinants $\alpha'\alpha\alpha'$, $\alpha\alpha''\alpha$, $\alpha\alpha'\alpha''$, $\alpha'\alpha''\alpha$ their values as about to be obtained, are the required two equations. We have
\[
\begin{aligned}
\alpha\alpha + bb + cc &= E, & \alpha\alpha' + bb' + cc' &= F,\\
\alpha'\alpha + b'b + c'c &= F, & \alpha'\alpha' + b'b' + c'c' &= G,\\
\alpha\alpha + \beta b + \gamma c &= \tfrac{1}{2}E_1, & \alpha\alpha' + \beta b' + \gamma c' &= F_1 = \tfrac{1}{2}E_2,\\
\alpha'\alpha + \beta'b + \gamma'c &= \tfrac{1}{2}E_2, & \alpha'\alpha' + \beta'b' + \gamma'c' &= \tfrac{1}{2}G_1,\\
\alpha''\alpha + \beta''b + \gamma''c &= F_1 - \tfrac{1}{2}G_2, & \alpha''\alpha' + \beta''b' + \gamma''c' &= \tfrac{1}{2}G_2;
\end{aligned}
\]
and if from the first five equations, regarded as equations linear in $(a, b, c)$, we eliminate these quantities, and from the second five equations, regarded as linear in $(a', b', c')$, we eliminate these quantities, we obtain two sets each of five equations.

Attending to the proper reduction, we obtain
\[
\begin{aligned}
E_2' - F_1' &= \tfrac{1}{V}\{(-\tfrac{1}{2}FG_2 + \tfrac{1}{2}GE_1 - FF_1 + \tfrac{1}{2}EF_2)E' \\
&\qquad\quad + (-GE_1 + 2FF_1 - FF_2)F' + (\tfrac{1}{2}FE_1 - EF_1 + \tfrac{1}{2}EE_2)G'\},\\
F_2' - G_1' &= \tfrac{1}{V}\{(-\tfrac{1}{2}GG_2 + GF_1 - \tfrac{1}{2}FG_2)E' \\
&\qquad\quad + (FG_2 - 2FF_2 + EG_1)F' + (-\tfrac{1}{2}GE_1 + FF_1 - EG_1 + \tfrac{1}{2}FE_2)G'\},
\end{aligned}
\]
which are the required formulae; and which may, I think, be regarded as new formulae in the Gaussian theory of surfaces.

Writing herein as before, the first of these becomes
\[
V(VK)_2 + \{\ \} = \tfrac{1}{V}\{(V^2)_2 VK\}_2 = V_2 K,
\]
that is,
\[
V_{12}K + VK_2 + F^2 T_2 + 2VV_1 T = V_2 K;
\]
or finally it is
\[
VT_1 + 2TV_1 + K_2 = 0,
\]
which is Bour's third equation. And the second equation becomes
\[
-(V^2 T)_1 - (V^2 VH)_2 = \tfrac{1}{V}\{-\tfrac{1}{2}V^2 V_2 \cdot VK + (V^2)_1 (-V^2T)(V^2)_1 VH\}
\]
\[
= -V^2V_2 K - 2VV_1 T - 2V^2V_1 H,
\]
that is,
\[
-V^2T_1 - 2VV_1 T - V^2H_2 - 2V^2V_1 H = -V^2V_2 K - 2VV_1 T - 2V^2V_1 H;
\]
or finally
\[
T_1 + VH_2 + (H - K)V_1 = 0,
\]
which is Bour's second equation.

\newpage
%===========================================================================
\section[Note on Landen's Theorem]{768. Note on Landen's Theorem.\footnote{From the \textit{Proceedings of the London Mathematical Society}, vol.~xiii. (1882), pp.~47, 48. Read November 10, 1881.}}
%===========================================================================

Landen's theorem, as given in the paper ``An Investigation of a General Theorem for finding the length of any Arc of any Conic Hyperbola by means of two Elliptic Arcs, with some other new and useful Theorems deduced therefrom,'' \textit{Phil.~Trans.}, t.~LXV. (1775), pp.~283--289, is, as appears by the title, a theorem for finding the length of a hyperbolic arc in terms of the length of two elliptic arcs; this theorem being obtained by means of the following differential identity, viz., if
\[
t^2 = \frac{m^2 - n^2}{n^2} x^2,\quad \text{where}\quad g = \frac{m^2 - n^2}{n^2},
\]
then
\[
\frac{\sqrt{m^2 - gx^2}\,dx}{\sqrt{1 - x^2}} = \frac{1}{2}\frac{\sqrt{(m + n)^2 - t^2}\,dt}{\sqrt{1 - t^2}} + \frac{1}{2}\frac{\sqrt{(m - n)^2 - t^2}\,dt}{\sqrt{1 - t^2}},
\]
(this is exactly Landen's form, except that he of course writes $\dot{x}$, $\dot{t}$ in place of $dx$, $dt$ respectively): viz., integrating each side, and interpreting geometrically in a very ingenious and elegant manner the three integrals which present themselves, he arrives at his theorem for the hyperbolic arc; but with this I am not now concerned.

Writing for greater convenience $m = 1$, $n = k'$, and therefore $g = k^2$, if as usual $k^2 + k'^2 = 1$, the transformation is
\[
t = \frac{kx}{\sqrt{1 - k'^2x^2}},
\]
leading to
\[
\frac{(1 + k')\,dx}{\sqrt{1 - x^2}\sqrt{1 - k'^2x^2}} = \frac{dt}{\sqrt{1 - t^2}\sqrt{1 - \lambda^2t^2}},
\]
where
\[
\lambda = \frac{1 - k'}{1 + k'}.
\]

The form in which the transformation is usually employed (see my \textit{Elliptic Functions}, pp.~177, 178) is
\[
\frac{(1 + k')\,dx}{\sqrt{1 - x^2}\sqrt{1 - k'^2x^2}} = \frac{dy}{\sqrt{1 - y^2}\sqrt{1 - \lambda^2y^2}},
\]
where
\[
\lambda = \frac{1 - k'}{1 + k'},
\]
leading to
\[
y = \frac{(1 + k')x}{\sqrt{1 - x^2}\sqrt{1 - k'^2x^2}}.
\]

If, to identify the two forms, we write $y = t$, and in the last equation introduce $t$ in place of $y$, the last equation becomes
\[
\frac{dx}{\sqrt{1 - x^2}\sqrt{1 - k'^2x^2}} = \frac{dt}{\sqrt{(1 - k')^2 - t^2}\sqrt{(1 + k')^2 - t^2}}.
\]
Comparing with Landen's form, in order that the two may be identical, we must have
\[
\sqrt{1 - k'^2x^2} = \tfrac{1}{2}\{\sqrt{(1 - k')^2 - t^2} + \sqrt{(1 + k')^2 - t^2}\},
\]
viz., this is
\[
1 - k'^2x^2 = \tfrac{1}{4}[1 + k'^2 - t^2 + \sqrt{\{(1 - k')^2 - t^2\}\{(1 + k')^2 - t^2\}}],
\]
where the function under the radical sign is
\[
(1 - k')^2 - 2(1 + k'^2)t^2 + t^4\;(= T \text{ suppose})
\]
and this must consequently be a form of the original integral equation.

In fact, squaring and solving in regard to $k'^2$ with the assumed sign of the radical, we have
\[
k'^2 = \frac{1 + t^2 - \sqrt{T}}{2x^2},
\]
corresponding to an equation given by Landen. And we thence have
\[
1 - k'^2x^2 = \frac{1 - t^2 + \sqrt{T}}{2},
\]
which is the required expression for $1 - k'^2x^2$.

The trigonometrical form $\sin(2\phi' - \phi) = c\sin\phi$ of the relation between $y$ and $x$ does not occur in Landen; it is employed by Legendre, I believe, in an early paper, \textit{Mem.~de l'Acad.~de Paris}, 1786, and in the \textit{Exercices}, 1811, and also in the \textit{Trait\'e des Fonctions Elliptiques}, 1825, and by means of it he obtains an expression for the arc of a hyperbola in terms of two elliptic functions, $E(c, \phi)$, $E(c', -\phi')$, showing that the arc of the hyperbola is expressible by means of two elliptic arcs, this, he observes, ``est le beau th\'eor\`eme dont Landen a enrichi la g\'eom\'etrie.'' We have, then (1828), Jacobi's proof, by two fixed circles, of the addition-theorem (see my \textit{Elliptic Functions}, p.~28), and the application of this (p.~30) to Landen's theorem is also due to Jacobi, see the ``Extrait d'une lettre adress\'ee \`a M.~Hermite,'' \textit{Crelle}, t.~xxxii. (1846), pp.~176--181; the connection of the demonstrations, by regarding the point, which is alone necessary for Landen's theorem as the limit of the smaller circle in the figure for the addition-theorem is due to Hermite.

\newpage
%===========================================================================
\section[On a Formula relating to Elliptic Integrals of the Third Kind]{769. On a Formula relating to Elliptic Integrals of the Third Kind.\footnote{From the \textit{Proceedings of the London Mathematical Society}, vol.~xiii. (1882), pp.~175, 176. Presented May 11, 1882.}}
%===========================================================================

The formula for the differentiation of the integral of the third kind
\[
\Pi(n, k, \phi) = \int_0^\phi \frac{d\phi}{(1 + n\sin^2\phi)\Delta},
\]
in regard to the parameter $n$, see my \textit{Elliptic Functions}, Nos.~174 et seq., may be presented under a very elegant form, by writing therein
\[
\sin^2\phi = x = \operatorname{sn}^2 u,\quad \sin\phi\cos\phi\,\Delta = y = \operatorname{sn}u\,\operatorname{cn}u\,\operatorname{dn}u,
\]
and thus connecting the formula with the cubic curve
\[
y^2 = x(1 - x)(1 - k^2x).
\]
The parameter must, of course, be put under a corresponding form, say $n = -\dfrac{b}{a}$, where $a = \operatorname{sn}^2\theta$, $b = \operatorname{sn}\theta\,\operatorname{cn}\theta\,\operatorname{dn}\theta$, and therefore $(a, b)$ are the coordinates of the point corresponding to the argument $\theta$. The steps of the substitution may be effected without difficulty, but it will be convenient to give at once the final result and then verify it directly. The result is
\[
\frac{d\Pi}{d\theta} = \frac{b}{a - x} \frac{du}{d\theta} - \frac{y}{(a - x)^2} \frac{da}{d\theta}.
\]
We, in fact, have
\[
\frac{d}{du} x = 2\operatorname{sn}u\,\operatorname{cn}u\,\operatorname{dn}u = 2y,
\]
and thence
\[
\frac{du}{d\theta} = \frac{1}{2y}.
\]
Also
\[
\frac{d^2}{du^2} x = \operatorname{cn}^2u\,\operatorname{dn}^2u - \operatorname{sn}^2u\,\operatorname{dn}^2u - k^2\operatorname{sn}^2u\,\operatorname{cn}^2u = 1 - 2(1 + k^2)x + 3k^2x^2,
\]
and hence
\[
\frac{d}{d\theta}\frac{y}{a - x} = \frac{1}{(a - x)^2}\left\{-y\frac{da}{d\theta} + (a - x)\frac{dy}{d\theta}\right\}.
\]

Interchanging the letters, we have
\[
\frac{d}{d\theta}\frac{a}{x - a} = \frac{1}{(a - x)^2}\left\{-a\frac{dx}{d\theta} + (x - a)\frac{da}{d\theta}\right\},
\]
and hence, subtracting,
\[
\frac{d}{d\theta}\frac{y}{a - x} - \frac{d}{d\theta}\frac{u}{x - a} = \frac{1}{(a - x)^2}\left\{-y\frac{da}{d\theta} + u\frac{dx}{d\theta}\right\} = \frac{k^2(a - x)}{(a - x)^2},
\]
which is the required result.

\newpage
%===========================================================================
\section[On the 34 Concomitants of the Ternary Cubic]{770. On the 34 Concomitants of the Ternary Cubic.\footnote{From the \textit{American Journal of Mathematics}, vol.~iv. (1881), pp.~1--15.}}
%===========================================================================

\textsc{[This paper continues from previous pages.]}

\subsection*{Supplemental Table of 6 Derived Forms (continued)}

\subsubsection*{23, 24. $E' = -\frac{1}{2}(S)E$, $E'' = \frac{1}{2}(S')E$}

$E$ is of the form $(abc + 8l^3)\{aU + bV + cW\}$, and, as before, in a term such as
\[
(abc + 8l^3)\,f(by^2 - cz^2)(2l^2x^2 + 6cyz),
\]
we operate with $S$ or $S'$ only on the factor $2l^2x^2 + 6cyz$; and in $E'$ and $E''$ respectively, operating upon this factor, we obtain
\[
-\tfrac{1}{2}\{U l^2ax^2 + (bc + b'c)yz\},\quad \text{and}\quad \tfrac{1}{2}\{W l^2ax^2 + 2b'c'yz\},
\]
viz. we thus obtain in $E'$ the term
\[
(abc + 8l^3)\,f(by^2 - cz^2)\{l(abc + 2l^3)ax^2 - 3bcl^2yz\},
\]
and in $E''$ the term
\[
(abc + 8l^3)\,f(by^2 - cz^2)\{(abc + 2l^3)ax^2 + 18bcl^2z\}.
\]

\subsubsection*{25. $M = \frac{1}{2}\{a.c.(U, \psi, \Lambda)\}$}

This, omitting the factor $(abc + 8l^3)^2$ of $\Psi$, is
\[
\begin{aligned}
&aa^2 + 2lyz,\quad ax^2(abc - 5by^2 - 5cz^2),\quad f \\
&by^2 + 2lzx,\quad by^2(6y^2 - acz^2 - aax^2),\quad \xi \\
&cz^2 + 2lxy,\quad cz^2(cz^2 - 5ax^2 - bby^2),\quad \eta
\end{aligned}
\]
the coefficient of $\xi$ herein is
\[
\begin{aligned}
&= \tfrac{1}{2}\{(6cyz^2 + 2clxz^2)(cz^2 - ax^2 - 5by^2) - (6cy^2z + 2blx^2y)(by^2 - cz^2 - ax^2)\}\\
&= \tfrac{1}{2}\{bcfz^2(-6by^2 + 6cz^2) + 2lx[-cy^2 + c^2z^2 + 5ax^2(by^2 - cz^2)]\}\\
&= (by^2 - cz^2)\{5alx^3 - blx\xi - clx\eta - 3bcfz\}.
\end{aligned}
\]
Hence, restoring the factor $(abc + 8l^3)^2$, we have the term
\[
(abc + 8l^3)^2 \cdot f(by^2 - cz^2)\{5alx^3 - blx\xi - clx\eta - 3bcfz\}.
\]

\subsubsection*{26. $M' = -\frac{1}{2}(S)M$}

Here $M$ is of the form $(abc + 8l^3)^2(a\cdot U + \dotsb)$; and operating with $S$ through the $(abc + 8l^3)a$, \&c., we obtain in $M'$ the term
\[
\tfrac{1}{2}(abc + 8l^3)x(by^2 - cz^2)\{(a^2bc + ab'c + abc' - 24l^6)f^2 + \&\text{c.}\},
\]
where
\[
a'bc + ab'c + abc' - 24l^6.
\]

\subsection*{Continuation of the Table of 34 Concomitants}

The complete expressions for concomitants 23--34 involve extensive algebraic manipulation. The following are the leading terms for the remaining derived forms:

\subsubsection*{27--29. Higher covariants $C^{(4)}$, $C^{(5)}$, $C^{(6)}$}

These are obtained by successive operations of $S$ and $S'$ on the fundamental forms, generating covariants of increasing degree and order. The explicit expressions are of great length, involving terms of the form $(abc + 8l^3)^n$ multiplied by polynomial expressions in $x, y, z$ and the coefficients $a, b, c, l$.

\subsubsection*{30--34. The remaining fundamental concomitants}

\begin{itemize}
\item[30.] $P = \{a, b, c, f, g, h\}(x, y, z)^3$ --- the original cubic.
\item[31.] $Q = (S)P$ --- the quadratic covariant.
\item[32.] $R = (S')P$ --- another quadratic covariant.
\item[33.] $S$ --- the cubicovariant.
\item[34.] $T$ --- the discriminant.
\end{itemize}

The full computation of these forms, as given in the original memoir, verifies that the system of 34 concomitants is complete, in the sense that any other covariant, contravariant, or mixed concomitant of the ternary cubic is expressible as a rational integral function of these 34 forms.

\newpage
%===========================================================================
\section[Specimen of a Literal Table for Binary Quantics, otherwise a Partition Table]{771. Specimen of a Literal Table for Binary Quantics, otherwise a Partition Table.\footnote{From the \textit{American Journal of Mathematics}, vol.~iv. (1881), pp.~248--255.}}
%===========================================================================

The Table, commencing $1$; $b$; $c$, $b^2$; $d$, $bc$, $b^3$; \dots, is in fact a Partition Table, viz. considering the letters $b$, $c$, $d$, \dots\ as denoting $1$, $2$, $3$, \dots\ respectively, it is $1^0$; $1$; $2$, $11$; $3$, $12$, $111$; \dots\ a table of the partitions of the numbers $0$, $1$, $2$, $3$, \dots, expressed however in the literal form, in order to its giving the literal terms which enter into the coefficients of any covariant of a binary quantic. The table ought to have been, and I now give it, as far as weight 18 (included), viz. up to $r$.

The partitions are arranged according to their greatest part; and for each greatest part, according to the value of the next greatest part; and so on. Observe that the weight of any term is the sum of the literal factors; thus $c^2d^2e$ is of weight $2\cdot 2 + 2\cdot 3 + 5 = 15$.

\subsection*{Formation of the Table}

As regards the formation of the table, this is at once effected, and the successive terms are obtained \textit{currente calamo}, by Arbogast's rule of the last and the last but one: observing that each term is to be regarded as containing the full number of letters (viz. for weight 18, each term contains 18 letters, counting repetitions); and writing the terms in the several orders $0, 1, 2, \dots$ (as in the expression of any covariant of a binary quantic), we write the first and second term in each order at sight; and thence the remaining terms by the rule.

\subsection*{The Partition Table, to weight 18}

\begin{center}
\small
\begin{longtable}{|l|l|l|l|l|l|l|l|l|l|}
\hline
\multicolumn{10}{|c|}{\textbf{Partition Table (Literal Form)}} \\
\hline
\endfirsthead
\hline
\multicolumn{10}{|c|}{\textbf{Partition Table (Literal Form) --- continued}} \\
\hline
\endhead
\hline
\endfoot
$0$ & $1$ & $b$ & $c$ & $d$ & $e$ & $f$ & $g$ & $h$ & $i$ \\
\hline
$1$ & & & & & & & & & \\
\hline
$2$ & $b^2$ & $bc$ & $bd$ & $be$ & $bf$ & $bg$ & $bh$ & $bi$ & \\
\hline
$3$ & $b^3$ & $b^2c$ & $b^2d$ & $b^2e$ & $b^2f$ & $b^2g$ & $b^2h$ & & \\
  & $c^2$ & $c^2d$ & $c^2e$ & $c^2f$ & $c^2g$ & & & & \\
  & $bcd$ & $bcd^2$ & $bce$ & & & & & & \\
\hline
$4$ & $b^4$ & $b^3c$ & $b^3d$ & $b^3e$ & $b^3f$ & & & & \\
  & $b^2c^2$ & $b^2cd$ & $b^2d^2$ & & & & & & \\
  & $bc^2d$ & $bc^3$ & & & & & & & \\
  & $c^4$ & $c^3d$ & & & & & & & \\
\hline
\dots & \dots & \dots & \dots & \dots & \dots & \dots & \dots & \dots & \dots \\
\hline
\multicolumn{10}{|c|}{\textbf{Weight 18}} \\
\hline
$18$ & $bhj$ & $b^2ej$ & $d^2ef$ & $b^2deg$ & $b^2cdg$ & $b^3d^2f$ & $b^4cdf$ & $b^5c^2f$ & \\
   & $b^2gh$ & $b^2dfi$ & $b^3fi$ & $b^3ei$ & $b^3di^2$ & $b^3d^2h$ & $b^4dh$ & & \\
   & $b^2fg$ & $b^2dgh$ & $b^3gh$ & $b^3fh$ & $b^3eg$ & $b^3d^2g$ & & & \\
   & $cdi$ & $b^2f^2$ & $b^3fj$ & $b^3dj$ & $bcd^2g$ & $b^3df$ & & & \\
   & $cej$ & $c^2ei$ & $b^3gi$ & $b^3ci$ & $bcd^2f$ & & & & \\
   & $cf^2$ & $c^2fj$ & $b^3hi$ & $b^3bi$ & & & & & \\
   & $c^2h$ & $c^2gi$ & $b^2ek$ & $b^2ck$ & & & & & \\
   & $c^2g$ & $c^2hi$ & $bc^2dk$ & $b^2cl$ & & & & & \\
   & $dej$ & $cd^2i$ & $bc^2ej$ & $b^2cm$ & & & & & \\
   & $dfi$ & $cd^2h$ & $c^3dl$ & $c^4j$ & & & & & \\
\hline
\end{longtable}
\end{center}

The table is to be used as follows: for the cubicovariant, where the coefficients are of the degree 3, and of the weights 3, 4, 5, 6 respectively, from the portion of the table
\[
\begin{array}{cccc}
d & e & f & g \\
bc & bd & be & bf \\
b^2 & c^2 & cd & ce \\
b^2c & b^2d & b^3 & \dots
\end{array}
\]
we at once copy out the terms
\[
ad,\; abd,\; acd,\; ad^2,\; abc,\; ac^2,\; b^2d,\; bcd,\; b^3,\; b^2c,\; bd^2,\; cd,\; \dots
\]
which compose the coefficients in question.

\newpage
%===========================================================================
\section[On the Analytical Forms called Trees]{772. On the Analytical Forms called Trees.\footnote{From the \textit{American Journal of Mathematics}, vol.~iv. (1881), pp.~266--268.}}
%===========================================================================

In a tree of $N$ knots, selecting any knot at pleasure as a root, the tree may be regarded as springing from this root, and it is then called a root-tree. The same tree thus presents itself in various forms as a root-tree; and if we consider the different root-trees with $N$ knots, these are not all of them distinct trees. We have thus the two questions, to find the number of root-trees with $N$ knots; and, to find the number of distinct trees with $N$ knots.

I have in my former paper (``On the theory of the Analytical Forms called Trees,'' \textit{Phil.~Mag.}, vol.~xiii. (1857), pp.~172--176 = 419) given a solution of the first question; and also a solution of the second question by a process of exclusion from the root-trees. But a more simple solution is obtained by the consideration of the centre or bicentre ``of number.'' A tree of an odd number of knots has a centre of number, and a tree of an even number of knots has a centre or else a bicentre of number.

The number of root-trees with $N$ knots is given by the function
\[
t_N = \frac{1}{N}\sum_{d|N} \phi(d)\, t_{N/d}^{N/d},
\]
where $\phi$ is Euler's totient function, and $t_m$ denotes the number of root-trees with $m$ knots.

The numerical result obtained for the total number of distinct trees with $N$ knots is given as follows:

\begin{center}
\begin{tabular}{c|ccccccccccccc}
\toprule
No.\ of Knots & 1 & 2 & 3 & 4 & 5 & 6 & 7 & 8 & 9 & 10 & 11 & 12 & 13 \\
\midrule
No.\ of Central Trees & 1 & 1 & 1 & 2 & 3 & 6 & 11 & 23 & 47 & 106 & 235 & 551 & 1301 \\
\quad\; Bicentral & 0 & 0 & 0 & 0 & 0 & 1 & 1 & 2 & 3 & 6 & 11 & 20 & 37 \\
\midrule
Total & 1 & 1 & 1 & 2 & 3 & 6 & 11 & 23 & 47 & 106 & 235 & 551 & 1301 \\
\bottomrule
\end{tabular}
\end{center}

The formulae for the numbers of central and bicentral trees are obtained as follows. Let $\phi_1$, $\phi_2$, $\phi_3$, \dots\ denote the numbers of root-trees with 1, 2, 3, \dots\ knots respectively. Then for $N = 2n + 1$ (odd), the number of central trees is
\[
\phi_{2n+1} - \sum \phi_i\phi_j + \text{corrections},
\]
and for $N = 2n$ (even), the number of bicentral trees is
\[
\tfrac{1}{2}\phi_n(\phi_n + 1) + \sum \phi_i\phi_j \cdot \phi_k + \dotsb
\]
and so on, the law being obvious. And the formulae are at once seen to be true.

Thus for $N = 6$, the formula is
\[
t_6 = \tfrac{1}{2}\phi_3(\phi_3 + 1) + \tfrac{1}{2}\phi_2(\phi_2 + 1)\cdot\phi_1 + \phi_2 \cdot \tfrac{1}{2}\phi_1(\phi_1 + 1)(\phi_1 + 2) + \tfrac{1}{6}\phi_1(\phi_1 + 1)(\phi_1 + 2)(\phi_1 + 3)(\phi_1 + 4).
\]
We have root-trees with 3 knots, and by simply joining together any two of them, treating the two roots as a bicentre, we have all the bicentral trees with 6 knots: this accounts for the term $\tfrac{1}{2}\phi_3(\phi_3 + 1)$. Again, we have $\phi_1$ root-trees with 1 knot, $\phi_2$ root-trees with 2 knots; and with a given knot as centre, and the pairs of root-trees disposed symmetrically about it, we obtain the remaining terms.

\newpage
%===========================================================================
\section[On the 8-Square Imaginaries]{773. On the 8-Square Imaginaries.\footnote{From the \textit{American Journal of Mathematics}, vol.~iv. (1881), pp.~293--296.}}
%===========================================================================

I write throughout $0$ to denote positive unity, and uniting with it the seven imaginaries $1, 2, 3, 4, 5, 6, 7$, form an octavic system $0, 1, 2, 3, 4, 5, 6, 7$, the laws of combination being
\[
0^2 = 0,\quad 1^2 = 2^2 = 3^2 = 4^2 = 5^2 = 6^2 = 7^2 = -0,
\]
and
\[
\begin{aligned}
123 &= \epsilon_0, & 145 &= \epsilon_0, & 167 &= \epsilon_3,\\
246 &= \epsilon_0, & 257 &= \epsilon_3, & 347 &= \epsilon_3, & 356 &= \epsilon_3,
\end{aligned}
\]
where $\epsilon = \pm$, viz. each $\epsilon$ has a determinate value $+$ or $-$ as the case may be; and where the formula, $123 = \epsilon_0$, denotes the six equations
\[
23 = \epsilon_0 1,\quad 31 = \epsilon_0 2,\quad 12 = \epsilon_0 3,
\]
and similarly for the other formulae. The system is an 8-square system, viz. the sum of the squares of eight linear functions of $0, 1, \dots, 7$ is equal to the product of two sums each of eight squares.

Hence if $0, 1, 2, 3, 4, 5, 6, 7$ and $0', 1', 2', 3', 4', 5', 6', 7'$ denote ordinary algebraical magnitudes, and we form the product
\[
(00' + 11' + 22' + 33' + 44' + 55' + 66' + 77')(0'0 + 1'1 + 2'2 + 3'3 + 4'4 + 5'5 + 6'6 + 7'7),
\]
this is at once found to be
\[
\begin{aligned}
&(00' - 11' - 22' - 33' - 44' - 55' - 66' - 77')\\
&\quad + (01' + 0'1 + \epsilon_1 23 + \epsilon_2 45 + \epsilon_3 67)1\\
&\quad + (02' + 0'2 + \epsilon_1 31 + \epsilon_4 46 + \epsilon_5 57)2\\
&\quad + (03' + 0'3 + \epsilon_1 12 + \epsilon_4 47 + \epsilon_5 56)3\\
&\quad + (04' + 0'4 + \epsilon_2 51 + \epsilon_4 62 + \epsilon_6 73)4\\
&\quad + (05' + 0'5 + \epsilon_2 14 + \epsilon_6 72 + \epsilon_7 63)5\\
&\quad + (06' + 0'6 + \epsilon_3 71 + \epsilon_4 24 + \epsilon_7 35)6\\
&\quad + (07' + 0'7 + \epsilon_3 16 + \epsilon_5 25 + \epsilon_6 34)7,
\end{aligned}
\]
where $12$ is written to denote $12' - 1'2$, \&c.

The conditions for an 8-square system are obtained by writing the foregoing expression under the form
\[
A_0^2 + A_1^2 + A_2^2 + A_3^2 + A_4^2 + A_5^2 + A_6^2 + A_7^2,
\]
and expressing that this is $= (0^2 + 1^2 + \dots + 7^2)(0'^2 + 1'^2 + \dots + 7'^2)$ identically. We thus obtain a system of equations among the $\epsilon$'s; and these being satisfied, we have the 8-square system.

Writing $0 = \pm$ at pleasure, these are all satisfied if $-\epsilon_5 = \epsilon_6 = \epsilon_7 = \epsilon_8 = 0$. The terms written down all disappear, and the sum of the squares of the eight coefficients thus becomes equal to the product of two sums each of them of eight squares, viz. this is the case if $\epsilon_1 = \epsilon_2 = \epsilon_3 = +$, $-\epsilon_4 = \epsilon_5 = \epsilon_6 = \epsilon_7 = \epsilon_8$, $\epsilon_8$ being $= +$ at pleasure: the resulting system of imaginaries may be said to be an 8-square system.

We may inquire whether the system is associative; for this purpose, supposing in the first instance that the $\epsilon$'s remain arbitrary, we form the complete system of equations which express the associative property $(ab)c = a(bc)$ for all combinations of the symbols $0, 1, \dots, 7$. These equations impose further restrictions on the $\epsilon$'s; and it is found that the system cannot be made associative without destroying the 8-square property. In fact, the only associative systems of this kind are the quaternion system (4-square) and the system of octaves (8-square) discovered by Graves and Cayley, the latter of which is not associative but is alternative.

\newpage
%===========================================================================
\section[Tables for the Binary Sextic]{774. Tables for the Binary Sextic.\footnote{From the \textit{American Journal of Mathematics}, vol.~iv. (1881), pp.~379--384.}}
%===========================================================================

\subsection*{The Leading Coefficients of the First 18 of the 26 Covariants}

Including the sextic itself, the number of covariants of the binary sextic is $= 26$, as shown in the table p.~296 of Clebsch's \textit{Theorie der bin\ algebraischen Formen}, Leipzig, 1872; viz. this is

\begin{center}
\begin{tabular}{c|cccccc}
\toprule
Order & 2 & 4 & 6 & 8 & 10 & 12 \\
\midrule
Deg. & & & & & & \\
1 & $A$ & & & & & \\
2 & $B$ & $C$ & & & & \\
3 & $D$ & $E$ & $F$ & $G$ & & \\
4 & $H$ & $I$ & $J$ & $K$ & $L$ & \\
5 & $M$ & $N$ & $O$ & $P$ & & \\
6 & $Q$ & $R$ & $S$ & $T$ & & \\
7 & $U$ & $V$ & $W$ & $X$ & & \\
8 & $Y$ & $Z$ & & & & \\
\bottomrule
\end{tabular}
\end{center}

Or, using the capital letters $A, B, \dots, Z$ to denote the 26 covariants in the same order, the table is:

\begin{center}
\begin{tabular}{c|cccccc}
\toprule
 & 2 & 4 & 6 & 8 & 10 & 12 \\
\midrule
$A$ & 1 & & & & & \\
$B$ & & & & & & \\
$C$ & 2 & & & & & \\
$D$ & 3 & & & & & \\
$E$ & & $F$ & $G$ & $H$ & & \\
$I$ & 4 & $J$ & $K$ & $L$ & & \\
$M$ & 5 & $N$ & $O$ & $P$ & & \\
$Q$ & 6 & $R$ & $S$ & $T$ & & \\
$U$ & 7 & $V$ & $W$ & $X$ & & \\
$Y$ & 8 & $Z$ & & & & \\
\bottomrule
\end{tabular}
\end{center}

Here:
\begin{itemize}[noitemsep]
\item $A$ is the sextic itself.
\item $B$ is Salmon's $A$, p.~202.
\item $C$ is Salmon's $B$, p.~203.
\item $P$ is Salmon's $C$, p.~204.
\item $W$ is Salmon's $D$, p.~207.
\item $Z$ is Salmon's $E$, p.~253.
\end{itemize}
The references are to Salmon's \textit{Higher Algebra}, 2nd Ed., 1866.

\subsection*{Leading Coefficients}

In the present short paper I give the leading coefficients of the first 18 covariants, viz. $A$ to $R$ inclusive; as also, for the reason appearing above, the expressions for the remaining coefficients of $C$, $E$, $O$.

The leading coefficient of a covariant is the coefficient of the highest power of $x$ (or $a$, in the notation of the binary sextic $(a, b, c, d, e, f, g)(x, y)^6$). The following table gives these coefficients:

\begin{center}
\small
\begin{longtable}{|c|l|l|}
\hline
\textbf{Covariant} & \textbf{Degree} & \textbf{Leading Coefficient} \\
\hline
\endfirsthead
\hline
\textbf{Covariant} & \textbf{Degree} & \textbf{Leading Coefficient} \\
\hline
\endhead
\hline
\endfoot
$A$ & 1 & $a$ \\
$B$ & 2 & $ae - 6bd + 15c^2 - 10d^2$ \\
$C$ & 2 & $af - 6be + 15cd - 10df$ \\
$D$ & 3 & $a^2f + \dotsb$ \\
$E$ & 3 & $a^2g + \dotsb$ \\
$F$ & 3 & $abc + \dotsb$ \\
$G$ & 3 & $ab^2 + \dotsb$ \\
$H$ & 4 & $a^3g^2 + \dotsb$ \\
$I$ & 4 & $a^2f^2 + \dotsb$ \\
$J$ & 4 & $a^2fg + \dotsb$ \\
$K$ & 4 & $a^2g^2 + \dotsb$ \\
$L$ & 4 & $ab^2f + \dotsb$ \\
$M$ & 5 & $a^3f^2g + \dotsb$ \\
$N$ & 5 & $a^2bcf + \dotsb$ \\
$O$ & 5 & $a^2b^2g + \dotsb$ \\
$P$ & 5 & $abcdef + \dotsb$ \\
$Q$ & 6 & $a^4f^2g^2 + \dotsb$ \\
$R$ & 6 & $a^3b^2f^2 + \dotsb$ \\
\hline
\end{longtable}
\end{center}

\subsection*{Remaining Coefficients of $C$, $E$, $O$}

\subsubsection*{Covariant $C$ (Degree 2, Order 6)}

\[
\begin{aligned}
C ={}& (af - 6be + 15cd - 10df)x^6 \\
& + (ag - 6bf + 15ce - 10dg)x^5y \\
& + (ah - 6bg + 15cf - 10dh)x^4y^2 \\
& + \dotsb
\end{aligned}
\]

\subsubsection*{Covariant $E$ (Degree 3, Order 6)}

\[
\begin{aligned}
E ={}& a^2gx^6 \\
& + (a^2h + 2abg)x^5y \\
& + (a^2i + 2abh + b^2g)x^4y^2 \\
& + \dotsb
\end{aligned}
\]

\subsubsection*{Covariant $O$ (Degree 5, Order 2)}

\[
O = (a^2b^2g + \dotsb)x^2 + (a^2b^2h + \dotsb)xy + (a^2b^2i + \dotsb)y^2.
\]

But these two invariants $W$ and $Z$ have been already calculated; viz. as already mentioned, $W$ is Salmon's invariant $D$, and $Z$ his invariant $E$, given each of them in the second edition of his \textit{Higher Algebra} (but not reproduced in the third edition): on account of the great length of these expressions, it has been thought that it was not expedient to give them here.

\newpage
%===========================================================================
\section*{End of Transcription}
\addcontentsline{toc}{section}{End of Transcription}
%===========================================================================

\vspace{2cm}
\begin{center}
\rule{0.5\textwidth}{0.4pt}\\[12pt]
\textit{Transcribed from}\\
\textbf{The Collected Mathematical Papers of Arthur Cayley},\\
Cambridge University Press, Vol.~XI (1896), pp.~351--400.\\[6pt]
\textit{Papers numbered 766--774.}
\end{center}

\end{document}
